\documentclass[../main]{subfiles}
\begin{document}
% ============================================================
% Front Matter: Title page, Abstract, Acknowledgements
% ============================================================

\begin{titlepage}
  \centering
  \vspace*{2cm}

  {\LARGE Kyoto University}\\[1cm]
  {\large Doctoral Thesis}\\[3cm]

  {\Huge\bfseries Microscopic Analysis of Fabric Evolution\\[0.5em]
   in Liquefaction Using Discrete Element Method}\\[3cm]

  {\Large HAN Yusong}\\[2cm]

  {\large Graduate School of Engineering\\
   Kyoto University}\\[2cm]

  {\large 2026}

  \vfill
\end{titlepage}

\cleardoublepage

% ---------- Abstract ----------
\chapter*{Abstract}
\addcontentsline{toc}{chapter}{Abstract}

Liquefaction, characterized by the loss of soil strength due to the buildup of pore water pressure during seismic events, poses significant risks to infrastructure and buildings. The Discrete Element Method (DEM) facilitates the direct analysis of particle interactions and microscopic fabric evolution, bridging the gap between micro-scale processes and macro-scale behaviors like liquefaction. This research establishes a methodology for evaluating seismic stability and liquefaction resistance using microscopic parameters with DEM.

Evolution of soil fabric under diverse stress states through drained true triaxial shear tests is explored using DEM. The distribution of contact density ($\rhoc$) evolves with increasing axial strain and different Lode angle during triaxial shear. The effects of stress anisotropy on liquefaction resistance of sand soils are investigated through undrained cyclic shear simulations.

A coupled analytical servo mechanism is developed to replicate undrained conditions and stress states in hollow cylinder apparatus (HCA) shear tests, in which the servo coefficients are derived from a volume-constrained stiffness matrix that accounts for the cross-coupling between radial and axial wall movements. The results demonstrate that the signed fabric anisotropy indicator correlates directly with the variation in liquefaction resistance among different \Kzero{} states, while the mechanical coordination number primarily reflects packing compactness. Experimental studies using HCA validate the non-monotonic trend in liquefaction resistance with stress anisotropy.

A new double-8 shear loading method is proposed to isolate and analyse the impact of multi-directional shear stress on liquefaction. Microscopic fabric metrics, such as \Zm{} and \Fc{}, play a vital role in evaluating liquefaction resistance. The variations in \Zm{} and fabric anisotropy during the liquefaction process reflect the susceptibility of granular sand to different shear stress patterns.

\vspace{1em}
\noindent\textbf{Keywords:} Liquefaction, Discrete element method, Fabric anisotropy, Hollow cylinder apparatus, Multi-directional shear stress, Coordination number

\cleardoublepage

% ---------- Acknowledgements ----------
\chapter*{Acknowledgements}
\addcontentsline{toc}{chapter}{Acknowledgements}

% TODO: Write acknowledgements

\cleardoublepage

\end{document}
